%% PLDI 2027 submission for strong-rep-nu.
%%
%% Class options follow the PLDI call for papers (PLDI 2026 wording; the
%% PLDI 2027 call was not yet posted on 2026-09-25): the PACMPL `acmsmall'
%% format, with review, screen, anonymous and nonacm.  At most 20 pages of
%% text, excluding the bibliography.  Double-blind: no author names, and
%% refer to our own prior work in the third person.
%%
%% The content plan is draft.md; the ACM sample this file is based on is
%% acmart-v2.20/acmsmall-conf.tex.
\documentclass[acmsmall,screen,review,anonymous,nonacm]{acmart}

\usepackage{semantic} % defines \inference
\usepackage{stmaryrd} % defines \fatsemi

\newcommand{\id}[1]{\key{id}(#1)}
\newcommand{\dual}[1]{-#1}
%% \newcommand{\seal}[1]{\key{seal}(#1)}
%% \newcommand{\unseal}[1]{\key{unseal}(#1)}
%% \newcommand{\seal}[1]{#1 \sharp}
%% \newcommand{\unseal}[1]{#1 \flat}
\newcommand{\seal}[1]{-#1}
\newcommand{\unseal}[1]{+#1}
\newcommand{\becomes}{\longrightarrow}
\newcommand{\of}{{:}}
\newcommand{\tu}{{\to}}
\newcommand{\lam}[1]{\lambda #1 \dotspace}
\newcommand{\Lam}[1]{\Lambda #1 \dotspace}
\newcommand{\all}[1]{\forall #1 \dotspace}
\newcommand{\app}{\;}
\newcommand{\Let}{\key{let}\:}
\newcommand{\In}{\:\key{in}\:}
\newcommand{\Int}{\key{Int}}
\newcommand{\Bool}{\key{Bool}}
\newcommand{\True}{\key{true}}
\newcommand{\False}{\key{false}}
\newcommand{\by}{{:=}}
\newcommand{\src}[1]{\meta{src}(#1)}
\newcommand{\trg}[1]{\meta{trg}(#1)}
\newcommand{\dom}{\key{dom}}
\newcommand{\rng}{\key{rng}}
\newcommand{\key}[1]{\mathtt{#1}}
\newcommand{\op}{\mathit{op}}
\newcommand{\dotspace}{.\,}
\newcommand{\comma}{,\,}
\newcommand{\sem}[1]{\llbracket #1 \rrbracket}
\newcommand{\nuinst}[4]{\nu #1\by #2 \dotspace (#3 \app #1) \langle #4 \rangle}
\newcommand{\boundry}[3]{\left[#2\right] #1 \langle #3 \rangle}

\newcommand{\black}[1]{{\color{black} #1}}
\newcommand{\gray}[1]{{\color{lightgray} #1}}
\definecolor{lightgray}{gray}{0.5}

\newcommand{\ba}{\begin{array}}
\newcommand{\ea}{\end{array}}
\newenvironment{stack}{\ba{@{}l@{}}}{\ea}
\newenvironment{branch}{\left\{\ba{@{}l@{\qquad}l@{}}}{\ea\right\}}
\newenvironment{syntax}{\[\ba{l@{\qquad}lcl}}{\ea\]}
\newenvironment{example}{\[\ba{l@{\:}l}}{\ea\]}

\newcommand{\mkid}[1]{\mathsf{Id}(#1)}
\newcommand{\isid}[1]{\mathsf{Id?}(#1)}
\newcommand{\cmp}{\mathbin{\fatsemi}}
\newcommand{\cmpt}{\mathbin{\fatsemi}}
\newcommand{\cmptt}{\mathbin{\fatsemi}}
\newcommand{\cmpg}{\mathbin{\fatsemi}}
\newcommand{\sm}{\mathbin{\fatsemi^{s}}}
\newcommand{\cancel}[2]{+#1(#2)}

\newcommand{\labeltag}[1]{\label{#1}\tag{\textsc{#1}}}
\newcommand{\tagsc}[1]{\tag{\textsc{#1}}}

\newcommand{\figtypes}{
\begin{figure*}[t]
\begin{flushleft}
Types
\begin{syntax}
  \text{Base types} & \iota & ::= & \Int \mid \Bool  \\
  \text{Type variables} & X,Y,Z \\
  \text{Rep. variables} & \alpha,\beta \\
  \text{Types} & A,B,C & ::= & X \mid \iota \mid A \to B \mid \all{X} A \\
  \text{Rep. Types} & R,S & ::= &
      \alpha \mid X \mid \iota  \mid R \to S \mid \all{X} R \\
  \text{Type Contexts} & \Delta & ::= &
    [] \mid \Delta, \alpha \mid \Delta, \alpha \by R \mid \Delta, X \by \alpha
\end{syntax}

Well-formed types \hfill \fbox{$\Delta \vdash A$}
\begin{gather*}
  \inference
   {X \by \alpha \in \Delta}
   {\Delta \vdash X}
  \quad
  \inference
   {}
   {\Delta \vdash \iota}
  \quad
  \inference
   {\Delta \vdash A & \Delta \vdash B}
   {\Delta \vdash A \to B}
  \quad
  \inference
   {\Delta, \alpha, X \by \alpha \vdash A}
   {\Delta \vdash \all{X} A} (X,\alpha \notin \Delta)
\end{gather*}
\end{flushleft}

\hfill \fbox{$X \by A \in \Delta$}
\[
X \by A \in \Delta \;\triangleq \;
  X \by \alpha \in \Delta
  \text{ and } \alpha \by R \in \Delta
  \text{ and } \Delta(A) = R
\]

\caption{Types and Type Contexts}
\label{fig:types}
\end{figure*}
}

\newcommand{\figctx}{
\begin{figure*}[t]
\begin{flushleft}
\begin{minipage}[t]{0.30\textwidth}
Scope changes
\[\ba{r@{\;}c@{\;}l}
  \delta & ::= &
    [] \mid \delta,{+X}^\alpha \mid \delta,-X^\alpha
\ea\]
\end{minipage}
\begin{minipage}[t]{0.35\textwidth}
\fbox{$\delta(\Delta)$} \quad interior
\begin{align*}
  [](\Delta) &= \Delta \\
  (\delta,{+X}^\alpha)(\Delta) &= \delta(\Delta),X\by \alpha \\
  (\delta,{-X}^\alpha)(\Delta) &= -X(\delta(\Delta))
\end{align*}
\end{minipage}
\begin{minipage}[t]{0.29\textwidth}
\fbox{$\dual{\delta}$} \quad inverse
\begin{align*}
  \dual{[]} &= [] \\
  \dual{(\delta,{+X}^\alpha)} &= [{-X}^\alpha]\mathop{++}\dual{\delta}\\
  \dual{(\delta,{-X}^\alpha)} &= [{+X}^\alpha]\mathop{++}\dual{\delta}
\end{align*}
\end{minipage}

\vspace{5pt}
\begin{minipage}[t]{0.44\textwidth}
\fbox{$-X(\Delta)$} remove $X$
\begin{align*}
  -X([]) &= [] \\
  -X(\Delta,X\by \alpha) &= \Delta \\
  -X(\Delta,Y\by \alpha) &= -X(\Delta),Y\by \alpha \quad \text{if } X \neq Y \\
  -X(\Delta,\alpha\by R) &= -X(\Delta),\alpha\by R \\
  -X(\Delta,\alpha) &= -X(\Delta),\alpha
\end{align*}
\end{minipage}
\begin{minipage}[t]{0.40\textwidth}
\fbox{$\Delta(A)$} \quad the representation type of $A$
\begin{align*}
  \Delta(X) &= \alpha \quad \text{if } X \by \alpha \in \Delta \\
  \Delta(X) &= X \quad \text{if } X \by \alpha \notin \Delta \text{ for any } \alpha \\
  \Delta(\iota) &= \iota \\
  \Delta(A \to B) &= \Delta(A) \to \Delta(B) \\
  \Delta(\all{X} A) &= \all{X} (-X(\Delta))(A)
\end{align*}
\end{minipage}


\vspace{6pt}
\noindent
\begin{minipage}[t]{0.62\textwidth}
\fbox{$\Delta \vdash \delta$} \quad coherent scope change
\begin{gather*}
  \inference{}{\Delta \vdash []}
  \quad
  \inference
   {\Delta \vdash \delta & X \by \alpha \notin \delta(\Delta) \\
    \forall\, Y \by \beta \in \delta^{+}(\Delta).\;
      X = Y \Leftrightarrow \alpha = \beta}
   {\Delta \vdash \delta,{+X}^\alpha}
  \\[2ex]
  \inference
   {\Delta \vdash \delta & X \by \alpha \in \delta(\Delta)}
   {\Delta \vdash \delta,{-X}^\alpha}
\end{gather*}
\end{minipage}%
\begin{minipage}[t]{0.36\textwidth}
\fbox{$\delta^{+}(\Delta)$} \quad conversion context
\begin{align*}
  []^{+}(\Delta) &= \Delta \\
  (\delta,{+X}^\alpha)^{+}(\Delta) &= \delta^{+}(\Delta) \cup \{X \by \alpha\} \\
  (\delta,{-X}^\alpha)^{+}(\Delta) &= \delta^{+}(\Delta)
\end{align*}
\end{minipage}
\end{flushleft}

\caption{Context Operations and Scope Changes.}
\label{fig:ctx}
\end{figure*}
}

\newcommand{\figconv}{
\begin{figure*}[t]
\begin{flushleft}
Conversions
\begin{syntax}
  \text{Conversion} & c,d & ::= & t \mid \unseal{X} \mid \unseal{X} ; c \\
  \text{Tail}   & t   & ::= & g \mid \seal{X} \mid t ; \seal{X} \\
  \text{Middle} & g,h & ::= & \id{A} \mid c \tu d \mid \all{X} c \\
  \text{Inert}  & i   & ::= & \id{X} \mid c \tu d \mid \all{X}c
    \mid \seal{X} \mid t; \seal{X} 
\end{syntax}

\noindent
\begin{minipage}[t]{0.30\textwidth}
\fbox{$\isid{c}$} \quad $c$ is an identity
\begin{align*}
  \isid{\id{A}} &= \mathit{true} \\
  \isid{c \tu d} &= \isid{c} \wedge \isid{d} \\
  \isid{\all{X} c} &= \isid{c} \\
  \isid{c} &= \mathit{false} \\
  & \quad \text{otherwise}
\end{align*}
\end{minipage}%
\begin{minipage}[t]{0.30\textwidth}
\fbox{$t \sm \seal{X}$, \ $\unseal{X} \sm d$} 
\begin{align*}
  t \sm \seal{X} &= \seal{X} \quad \text{if $\isid{t}$} \\
  t \sm \seal{X} &= t ; \seal{X} \quad \text{otherwise} \\
  \unseal{X} \sm t &= \unseal{X} \quad \text{if $\isid{t}$} \\
  \unseal{X} \sm t &= \cancel{X}{t} \quad \text{otherwise} \\
  \unseal{X} \sm d &= \unseal{X} ; d \\
  &\text{if $d$ begins with unseal}
\end{align*}
\end{minipage}%
\begin{minipage}[t]{0.30\textwidth}
\fbox{$\cancel{X}{t}$}
\begin{align*}
  \cancel{X}{g} &= \unseal{X} ; g \\
  \cancel{X}{\seal{X}} &= \id{X} \\
  \cancel{X}{\seal{Y}} &= \unseal{X} ; \seal{Y} \quad \text{if } X \neq Y\\
  \cancel{X}{t ; \seal{Y}} &= \cancel{X}{t} \sm \seal{Y} \quad \text{if $\cancel{X}{t}$ a tail} \\
  \cancel{X}{t ; \seal{Y}} &= \unseal{X} ; (t ; \seal{Y})
    \quad \text{otherwise}
\end{align*}
\end{minipage}

Well-typed conversions \hfill \fbox{$\Delta \mid \delta \vdash c : R \Rightarrow S$}
\begin{gather*}
  \inference
   {}
   {\Delta \mid \delta \vdash \id{\iota} : \iota \Rightarrow \iota}
  \quad
  \inference
   {X \by \alpha \in \Delta}
   {\Delta \mid \delta \vdash \id{X} : \alpha \Rightarrow \alpha}
  \\[2ex]
  \inference
   {\Delta \mid \dual{\delta} \vdash c : R' \Rightarrow R &
    \Delta \mid \delta \vdash d : S \Rightarrow S'}
   {\Delta \mid \delta \vdash c \tu d : R \to S \Rightarrow R' \to S'}
  \\[2ex]
  \inference
   {\Delta, \alpha, X \by \alpha \mid \delta \vdash c :
      R[X \mapsto \alpha] \Rightarrow S[X \mapsto \alpha]}
   {\Delta \mid \delta \vdash \all{X} c : \all{X} R \Rightarrow \all{X} S}
   (X,\alpha \notin \Delta)
  \\[2ex]
  \inference
   {\alpha \by \Delta(A) \in \Delta & {-X}^\alpha \in \delta}
   {\Delta \mid \delta \vdash \seal{X} : \Delta(A) \Rightarrow \alpha}
  \quad
  \inference
   {\Delta \mid \delta \vdash t : R \Rightarrow \Delta(A) \\
    \alpha \by \Delta(A) \in \Delta & {-X}^\alpha \in \delta &
    t \sm \seal{X} = t ; \seal{X}}
   {\Delta \mid \delta \vdash t ; \seal{X} : R \Rightarrow \alpha}
  \\[2ex]
  \inference
   {\alpha \by \Delta(A) \in \Delta & {+X}^\alpha \in \delta}
   {\Delta \mid \delta \vdash \unseal{X} : \alpha \Rightarrow \Delta(A)}
  \quad
  \inference
   {\alpha \by \Delta(A) \in \Delta & {+X}^\alpha \in \delta &
    \unseal{X} \sm c = \unseal{X} ; c \\
    \Delta \mid \delta \vdash c : \Delta(A) \Rightarrow S}
   {\Delta \mid \delta \vdash \unseal{X} ; c : \alpha \Rightarrow S}
\end{gather*}
\end{flushleft}

\caption{Conversions}
\label{fig:conv}
\end{figure*}
}

\newcommand{\fignufstatics}{
\begin{figure*}[t]
\begin{flushleft}
Terms
\begin{syntax}
  \text{Variables} & x,y  \\
  \text{Constants} & k &::=&
    \key{0} \mid \key{1} \mid \cdots \mid \True \mid \False \mid \cdots \\
  \text{Operators} & \op &::=& \key{+} \mid \key{\leq} \mid \cdots \\
  \text{Terms} & L,M,N & ::= &
  k \mid \op(\vec{M}) \mid x \mid \lam{x \of A}N \mid L \app M  \\
  & & & \Lam{X} V \mid \nuinst{X}{A}{L}{c} \mid \boundry{M}{\delta}{c} \\
  \text{Simples} & U & ::= &
     k \mid \lam{x \of A}N \mid \Lam{X}V \\
  \text{Values} & V,W & ::= & U \mid \boundry{U}{\delta}{i}
\end{syntax}

\begin{syntax}
  \text{Type Environment} & \Gamma & ::= &
    [] \mid \Gamma \comma x:A 
\end{syntax}
Term typing \hfill \fbox{$\Delta\mid\Gamma \vdash M:A$}
\begin{gather*}
  \gray{\inference
   {k:\iota}
   {\Delta\mid\Gamma \vdash k : \iota}}
  \quad
  \gray{\inference
   {\Delta\mid\Gamma \vdash \vec{M} : \vec{\iota} & \op : \vec{\iota} \to \iota}
   {\Delta\mid\Gamma \vdash \op(\vec{M}\,) : \iota}}
  \quad
  \gray{\inference
   {x:A \in \Gamma}
   {\Delta\mid\Gamma \vdash x : A}}
  \\[2ex]
  \gray{\inference
   {\Delta\mid\Gamma,x:A \vdash N : B}
   {\Delta\mid\Gamma\vdash \lam{x \of A} N : A \tu B}}
  \quad
  \gray{\inference
   {\Delta\mid\Gamma \vdash L : A \tu B & \Delta\mid\Gamma \vdash M : A }
   {\Delta\mid\Gamma \vdash L \app M : B}}
  \quad
  \gray{\inference
   {\Delta,\alpha,X \by \alpha \mid \Gamma \vdash V : A}
   {\Delta \mid \Gamma \vdash \Lam{X}V : \all{X}A}}
  \\[2ex]
  \inference
   {\Delta \vdash A & \Delta \mid \Gamma \vdash L : \all{X} C \\
    \Delta' \mid [{+X}^\alpha] \vdash c : \Delta'(C) \Rightarrow \Delta(B) &
    \Delta \vdash B}
   {\Delta \mid \Gamma \vdash \nuinst{X}{A}{L}{c} : B}
   (\Delta' = \Delta,\alpha\by\Delta(A),X\by \alpha)
  \\[2ex]
  %BoundaryWf \Delta Θ \Delta_i \Delta_c
  \inference
  {\Delta \vdash \delta &
   \delta(\Delta) \mid [] \vdash M : A \\
   \delta^{+}(\Delta) \mid \delta \vdash c :
     \delta(\Delta)(A) \Rightarrow \Delta(B) &
   \Delta \vdash B
  }
  {\Delta \mid \Gamma \vdash \boundry{M}{\delta}{c} : B}
\end{gather*}

\end{flushleft}

\caption{System $\nu$F: Static Semantics}
\label{fig:nu-F}
\end{figure*}
}

\newcommand{\figcompose}{
\begin{figure*}[t]
\begin{flushleft}
\small
\noindent
\begin{minipage}[t]{0.45\textwidth}
\fbox{$\Delta \vdash c \cmp d$} \quad compose conversions
\begin{align*}
  \Delta \vdash \unseal{X} \cmp d &= \Delta \vdash \unseal{X} \sm d \\
  \Delta \vdash (\unseal{X} ; c) \cmp d &= \unseal{X} \sm (\Delta \vdash c \cmp d) \\
  \Delta \vdash t \cmp d &= \Delta \vdash t \cmpt d
\end{align*}
\fbox{$\Delta \vdash t \cmpt d$} \quad a tail, then a conversion
\begin{align*}
  \Delta \vdash t \cmpt t' &= t \cmptt t' \\
  \Delta \vdash \seal{X} \cmpt \unseal{Y} &= \mkid{A} 
    \;\; (X \by A \in \Delta) \\
  \Delta \vdash \seal{X} \cmpt (\unseal{Y} ; c) &= c \\
  \Delta \vdash (t ; \seal{X}) \cmpt \unseal{Y} &= t \\
  \Delta \vdash (t ; \seal{X}) \cmpt (\unseal{X} ; d) &= \Delta \vdash t \cmpt d \\
  \Delta \vdash \id{Y} \cmpt \unseal{Y} &= \unseal{Y} \\
  \Delta \vdash \id{Y} \cmpt (\unseal{Y} ; d) &= \unseal{Y} ; d
\end{align*}
\end{minipage}
\begin{minipage}[t]{0.45\textwidth}
\fbox{$\Delta \vdash t \cmptt t'$} \quad a tail, then a tail
\begin{align*}
  \Delta \vdash t \cmptt \seal{Y} &= t \sm \seal{Y} \\
  \Delta \vdash t \cmptt (t' ; \seal{Y}) &= (\Delta \vdash t \cmptt t') \sm \seal{Y} \\
  \Delta \vdash g \cmptt h &= \Delta \vdash g \cmpg h \\
  \Delta \vdash \seal{X} \cmptt \id{X} &= \seal{X} \\
  \Delta \vdash (t ; \seal{X}) \cmptt \id{X} &= t ; \seal{X}
\end{align*}
\fbox{$\Delta \vdash g \cmpg h$} \quad a middle, then a middle
\begin{align*}
  \Delta \vdash \id{A} \cmpg h &= h \\
  \Delta \vdash (c \tu d) \cmpg \id{B} &= c \tu d \\
  \Delta \vdash (c \tu d) \cmpg (c' \tu d') &= (\Delta \vdash c' \cmp c) \tu (\Delta \vdash d \cmp d') \\
  \Delta \vdash \all{X} c \cmpg \id{B} &= \all{X} c \\
  \Delta \vdash \all{X} c \cmpg \all{X} c' &=
      \all{X} (\Delta, \alpha, X \by \alpha \vdash c \cmp c')
\end{align*}
\fbox{$\mkid{A}$} \quad the identity at $A$
\begin{align*}
  \mkid{\iota} &= \id{\iota} \\
  \mkid{X} &= \id{X} \\
  \mkid{A \to B} &= \mkid{A} \tu \mkid{B} \\
  \mkid{\all{X} A} &= \all{X} \mkid{A}
\end{align*}
\end{minipage}
\end{flushleft}

\caption{Composition of Conversions.}
\label{fig:compose}
\end{figure*}
}

\newcommand{\fignufdynamics}{
\begin{figure*}[t]
\begin{flushleft}

\begin{syntax}
  \text{Frames} & F & ::= & \Box \app M \mid V \app M \mid \nuinst{X}{A}{\Box}{c}
     \mid \boundry{\Box}{\delta}{c} \\
  \text{Rep. Alloc.} & \xi & ::= & \epsilon \mid \alpha := R 
\end{syntax}

\fbox{$F(\Delta)$} \quad the context at a frame's hole
\begin{align*}
  (\Box\app M)(\Delta) &= \Delta \\
  \cdots \\
  (\boundry{\Box}{\delta}{c})(\Delta) &= \delta(\Delta) 
\end{align*}

\vspace{10pt}
\fbox{$\Delta \vdash M \becomes N \dashv \xi $}
\begin{align}
  \Delta \vdash \op(\vec{V})
    & \becomes~ \sem{\op}(\vec{V}) \dashv \epsilon
    \label{Delta}\tagsc{Delta}  \\
  \Delta \vdash (\lam{x \of A} N) \app V
    & \becomes~ N[x \by V] \dashv \epsilon
    \label{Beta}\tagsc{Beta}  \\
  \Delta \vdash \boundry{U}{\delta}{c \tu d} \app W &
    \becomes \boundry{U \app (\boundry{W}{\dual{\delta}}{c})}{\delta}{d}
    \dashv \epsilon 
    \label{Wrap}\tagsc{Wrap} \\
  \Delta \vdash \nuinst{X}{A}{(\Lam{X}V)}{d} &\becomes~
          \boundry{V}{{+X}^\alpha}{d} \dashv \alpha\by \Delta(A)
          \label{TyBeta}\tagsc{TyBeta}\\
   \Delta \vdash \nuinst{X}{A}{\boundry{(\Lam{X}V)}{\delta}{\all{X}c}}{d} &
       \becomes~ \boundry{(\boundry{V}{\delta}{c})}{{+X}^\alpha}{d}
        \dashv \alpha\by \Delta(A)
        \label{TyWrap}\tagsc{TyWrap} \\
   \Delta \vdash  \boundry{\boundry{U}{\delta_1}{t_1}}{\delta_2}{c_1} &
     \becomes~ \boundry{U}{\delta_2 \mathop{++} \delta_1}{d}  \dashv \epsilon \label{Merge}\tagsc{Merge} \\
     & \text{where } (\delta_2 \mathop{++} \delta_1)^{+}(\Delta) \vdash t_1 \cmp c_1  = d
       \notag \\
   \Delta \vdash \boundry{U}{\delta}{\id{\iota}} &\becomes~ U \dashv \epsilon
     \label{Id}\tagsc{Id} \\
   \Delta \vdash F[M] &\becomes~ F[M'] \dashv \xi \qquad \text{if } F(\Delta) \vdash M \becomes M' \dashv \xi \label{Xi}\tagsc{$\xi$}
\end{align}

\noindent{\small With de Bruijn indices (the mechanization), the
conversions that \textsc{Wrap} and \textsc{Merge} move to another
boundary are renamed there, each type variable to the new position's
name for the same representation variable.  The renaming is computed
from the scopes alone, so \textsc{Wrap} reads no context and
\textsc{Merge} reads only $(\delta_2 \mathop{++} \delta_1)^{+}(\Delta)$.
With names, as here, no renaming is needed.}
\end{flushleft}

\caption{System $\nu$F: Dynamic Semantics}
\label{fig:dynamics}
\end{figure*}
}
  
\begin{document}

\title{System $\nu$F: Revisiting Syntactic Type Abstraction}  % working title

%% Hidden by the `anonymous' option in the submitted PDF.
\author{Anonymous Author(s)}
\affiliation{%
  \institution{Anonymous Institution}
  \country{Anonymous Country}}

\begin{abstract}
  % TODO
\end{abstract}

\maketitle

\section{Introduction}

\figtypes
\figctx
\figconv
\fignufstatics
\figcompose
\fignufdynamics

\section{Why a Function Domain Is Checked Under the Inverse Scope}
\label{sec:inverse-scope}

A boundary $\boundry{M}{\delta}{c}$ checks its conversion $c$ under its
scope change $\delta$ (Figure~\ref{fig:nu-F}), and $\delta$ grants
\emph{permissions}: a seal $\seal{X}$ may conceal the representation
variable $\alpha$ only if $\delta$ removes it from the interior's view
(an entry ${-X}^\alpha$), and an unseal $\unseal{X}$ may reveal $\alpha$
only if $\delta$ binds it there (an entry ${+X}^\alpha$).  The
permissions keep a conversion from sealing at an arbitrary type
variable that happens to be in scope: every seal and unseal must cite
the scope change that makes its variable abstract on one side of the
boundary and concrete on the other.  One rule of the conversion
judgement (Figure~\ref{fig:conv}) does not simply pass the scope down:
the domain of a function conversion is checked under the inverse
$\dual{\delta}$.  This section explains why, and why the inverse is the
right scope rather than an approximation of one.

\paragraph{The entry names the variable.}
The seal and unseal rules do not look $X$ up in the context: the entry
${-X}^\alpha$ or ${+X}^\alpha$ already links $X$ to $\alpha$.  This is
sound because the judgement is only ever entered as
$\delta^{+}(\Delta) \mid \delta \vdash c$ with $\Delta \vdash \delta$
(the boundary and $\nu$ rules), and a coherent $\delta$ puts the binding
of every entry into $\delta^{+}(\Delta)$: a bind ${+X}^\alpha$ adds
$X \by \alpha$ itself, and an unbind ${-X}^\alpha$ requires
$X \by \alpha$ in the interior at that point, which is contained in the
union.  So $X \by \alpha \in \delta^{+}(\Delta)$ is implied.  The fact
survives the rules that recurse: a domain checked under $\dual{\delta}$
keeps the context, and $\dual{\delta}$ has the same entries with the
signs flipped; the body of $\all{X} c$ only extends the context.  A
proof about conversion typing that needs the name lookup, for instance
that a conversion's types are unique or that composition is well
typed, can therefore simply include a premise that the scope is
coherent: the context is $\delta^{+}(\Delta)$ for some $\Delta$ with
$\Delta \vdash \delta$.  (The rule for $\id{X}$ keeps its lookup
$X \by \alpha \in \Delta$, since there is no entry to supply the name.)

\paragraph{The domain of a function conversion runs the other way.}
In $c \tu d$, the codomain conversion $d$ carries a result from the
interior of the boundary to its exterior, but the domain conversion $c$
carries an argument from the exterior to the interior.  So in $c$ the
roles of concealing and revealing are swapped.  The simplest
instantiation already shows this.  Instantiating the identity at
$\Int$,
\[
  \Delta \vdash \nuinst{X}{\Int}{(\Lam{X}\lam{x \of X} x)}{\seal{X} \tu \unseal{X}}
  \;\becomes\;
  \boundry{(\lam{x \of X} x)}{{+X}^\alpha}{\seal{X} \tu \unseal{X}}
  \dashv \alpha \by \Int ,
\]
gives a boundary whose scope has a single entry, ${+X}^\alpha$.  The
codomain $\unseal{X}$ reveals $\alpha$ to the exterior, and ${+X}^\alpha$
grants that.  The domain $\seal{X}$ conceals the incoming $\Int$ as $X$
on its way in, but the only entry that mentions $\alpha$ is
${+X}^\alpha$, not ${-X}^\alpha$.  If the domain were checked under
$\delta$ like the codomain, the most basic instantiation conversion
would be ill-typed.  Under the inverse it is fine: inverting a scope
flips the sign of every entry, so ${-X}^\alpha \in \dual{\delta}$
exactly when ${+X}^\alpha \in \delta$, and ${+X}^\alpha \in \dual{\delta}$
exactly when ${-X}^\alpha \in \delta$.  The domain seal needs the entry
${-X}^\alpha$ in $\dual{[{+X}^\alpha]} = [{-X}^\alpha]$, which is
there.  The type variables a conversion mentions are still
resolved in the same context $\delta^{+}(\Delta)$; only the source of
the permissions changes.

\paragraph{The inverse is exactly what \textsc{Wrap} needs.}
The rule
$\boundry{U}{\delta}{c \tu d} \app W \becomes
 \boundry{U \app (\boundry{W}{\dual{\delta}}{c})}{\delta}{d}$
(Figure~\ref{fig:dynamics}) moves the domain $c$ out of the function
conversion and makes it the conversion of a new boundary whose scope is
$\dual{\delta}$.  Before the step, $c$ was checked under $\dual{\delta}$
as the domain of $c \tu d$; after the step, it is checked under
$\dual{\delta}$ as the conversion of its own boundary.  Its typing
carries over unchanged, with the same representation types and the same
scope, so \textsc{Wrap} needs no argument about permissions at all.
\textsc{Merge} is almost as simple: $\delta_2 \mathop{++} \delta_1$
contains every entry of $\delta_1$ and of $\delta_2$, so both
conversions keep their permissions under the merged scope (inside a
domain, under its inverse, which likewise contains the inverses of
both), and they are composed under one scope.  In this reading, the
scope a subconversion is checked under is the scope of the boundary it
would become if a reduction pulled it out.

\paragraph{One scope per position, not one direction per conversion.}
One might instead give a whole conversion a single direction, recording
whether it conceals or reveals.  That does not work, because one
conversion may seal one type variable in its domain and another in its
codomain.  For instance
\[
  \seal{Y} \tu \seal{X}
  \quad\text{under}\quad
  \delta = [{+Y}^\beta, {-X}^\alpha]
\]
conceals the argument as $Y$ on its way in and the result as $X$ on its
way out, and no single direction describes both.  Checking the domain
under the inverse records nothing about the direction of the
conversion; the direction of each type variable comes from its own
entry in $\delta$.  The example is well typed: the domain seal
$\seal{Y}$ needs ${-Y}^\beta \in \dual{\delta}$, which holds because
${+Y}^\beta \in \delta$, and the codomain seal $\seal{X}$ needs
${-X}^\alpha \in \delta$.

\paragraph{Mechanization note.}  The first mechanization carried a
polarity bit on the judgement instead, reading $\delta$ at $+$ and
$\dual{\delta}$ at $-$.  The two formulations are interchangeable.
Checking under the inverse needs two small lemmas the bit avoids
(membership in an inverse scope, and the fact that renaming
representation variables commutes with inversion), and in exchange
it removes the polarities, their flipping lemmas, and the permission
map \textsc{Wrap} needed.  With de Bruijn indices the mechanization
checks only the representation variable of a permission (some entry
${-Y}^\alpha$ for a seal), because the index in a scope entry and the
index in the conversion live in different contexts and cannot be
compared; for the same reason its seal and unseal leaves keep the
lookup of $X$ in the conversion context, which is the only link from
the index to $\alpha$.  With names nothing is lost: the conversion
context names each representation variable at most once (the coherence
condition on ${+X}^\alpha$), and every entry's binding is in it, so an
entry for $\alpha$ always carries the name $X$ that the conversion
uses.

\bibliographystyle{ACM-Reference-Format}
\bibliography{references}

\end{document}
